\documentclass[11pt]{article}
\usepackage{amsmath,amssymb,amsthm}
\usepackage{mathtools}
\usepackage{tikz}
\usepackage{enumitem}
\usepackage[margin=1in]{geometry}

\newtheorem{theorem}{Theorem}
\newtheorem*{theorem*}{Theorem}
\newtheorem{lemma}{Lemma}
\newtheorem{definition}{Definition}
\newtheorem*{definition*}{Definition}
\newtheorem{proposition}{Proposition}
\newtheorem{assumption}{Assumption}
\newtheorem{remark}{Remark}



\title{Math 106 -- Notes \\ Week 7: February 18, 2026}
\author{Ethan Levien}
\date{}

\begin{document}
\maketitle


\section*{Forward and backward equations for It\^o Diffusion (Section 8.1/8.5)}
I will do the 1D case, but the calculations are easily generalized. We consider
\begin{equation}
dX_t = b(X_t)dt + \sigma(X_t)dW_t
\end{equation}
By It\^o's formula
\begin{equation}
f(X_t) = f(X_0) + \int_0^t(f'(X_s)b(X_s) + \frac{1}{2}f''(X_s)\sigma^2(X_s))ds + \int_0^tf'(X_s)\sigma(X_s)dW_s.
\end{equation}
Let
\begin{equation}
u(x,t) = {\mathbb E}[f(X_t)|X_0=x].
\end{equation}
Hence
\begin{align}
u(x,t) &= f(x) +  \int_0^t{\mathbb E}[f'(X_s)b(X_s) + \frac{1}{2}f''(X_s)\sigma^2(X_s)]ds\\
& \implies \partial_t u(x,t)\Big|_{t=0} = f'(x)b(x) + \frac{1}{2}f''(x)\sigma^2(x)
\end{align}
This means the generator is 
\begin{equation}
{\mathcal L} = b(x)\partial_x + \frac{1}{2}\sigma^2(x)\partial_x^2.
\end{equation}
This gives the backward equation
\begin{equation}
\partial_t u(x,t) = {\mathcal L} u(x,t),\quad u(x,0) = f(x).
\end{equation}
The forward equation is 
\begin{equation}
\partial_t p(x,t) = {\mathcal L}^* p(x,t)
\end{equation}
This is also called the \emph{Fokker-Planck} equation in the context of diffusions. 
By the definition of adjoint,
\begin{equation}
\langle  {\mathcal L}^*p,f \rangle = \langle p,  {\mathcal L}f \rangle
\end{equation}
and we therefore have
\begin{equation}
{\mathcal L}p(x) =-\partial_x[ b(x)p(x)] + \frac{1}{2}\partial_x^2[\sigma^2(x)p(x)].
\end{equation}
Often we consider the forward equation with initial conditions $p(x,t) = \delta(x-x_0)$, in which case the density does not exist at $0$, but we can still interpret the PDE in the weak sense. 

\section*{Example 8.1: Brownian Dynamics }
Consider a probability density $p$ which distributions the density of particles in a physical system. We will study a certain gradient flow on the space of probability densities. The starting point is the idea of \emph{free energy}, defined as 
\begin{equation}
F[p] = U[p] - T S[p]
\end{equation}
where
\begin{align}
U[p]  &= \int_{{\mathbb R}^d} U(x)dx\quad \text{(average of internal energy $U(x)$)}\\
S[p] &= k_B\int_{{\mathbb R}^d}p(x)\ln p(x)dx \quad \text{(entropy -- large means more uncertainty)}
\end{align}
A physical system will evolve to a distribution which minimizes $F[p]$, and these can be understood as distributions which minimize the energy while maximizing entropy. The temperature $T$ is a Lagrange multiplier. 
To do so we consider the variational derivative. For a perturbation $\eta$, 
\begin{equation}
\frac{d}{d\varepsilon}F[p + \epsilon \eta]\Big|_{\varepsilon =0} = \int_{{\mathbb R}^d} \left\{ U(x)\eta(x) + k_B T\frac{d}{d\varepsilon}(p + \varepsilon \eta)\ln (p + \varepsilon \eta) \right\} \Big|_{\varepsilon =0} dx
\end{equation}
and 
\begin{equation}
\frac{d}{d\varepsilon}(p + \varepsilon \eta)\ln (p + \varepsilon \eta) \Big|_{\varepsilon =0} = \eta + \eta \ln p
\end{equation}
Therefore, if we Taylor expand $F[p]$, 
\begin{equation}
F[p + \varepsilon \eta] = F[p] + \varepsilon \int_{{\mathbb R}^d} \eta(x) \mu(x)dx + O(\varepsilon^2)
\end{equation}
we have introduced the \emph{chemical potential}
\begin{equation}
\mu(x) = \frac{\delta F}{\delta p} = U(x) + k_B T(1 + \ln p(x)).
\end{equation}
The chemical potential tells us how the $O(\varepsilon)$ corrections to the free energy when we change $p$ by $\varepsilon$. For a system in an equilibrium the chemical potential is a constant which is determined by the normalization of $p(x)$. 
\begin{equation}
p(x) = \frac{1}{Z} e^{-U(x)/k_BT},\quad Z = \int_{\mathbb R^d}e^{-U(x)/k_BT}dx.
\end{equation}
For a probability distribution not in equilibrium, we can consider \emph{Fick's law} tells us that the rate at which a particle 
\begin{equation}
u(x) = - \frac{1}{\gamma} \nabla \mu(x)
\end{equation}
where $\gamma$ is a friction coefficient. 
This says that particles will flow along gradients in the chemical potential and the flow is faster with smaller frication. This assumes we are in the over-damped limit where these is no inertia, hence we do not need to treat the velocity as a separate dependence variable.   

The current (a probability per unit time) at point $x$ is 
\begin{equation}
j(x,t) = p(x,t)u(x).
\end{equation}
Conservation of mass means that in a region $A \subset {\mathbb R}^d$, 
\begin{equation}
\frac{d}{dt}\int_A p(x,t)dx = - \int_{\partial A} j(x) \cdot n(x)dx
\end{equation}
where $n(x)$ is the normal vector to the boundary $\partial A$ at  $x \in \partial A$. By the divergence Theorem, we have
\begin{equation}
 - \int_{\partial A} j(x) \cdot n(x)dx = \int_A \nabla \cdot j(x) dx
\end{equation}
and this holds for any $A$, hence
\begin{equation}
\partial_t p(x,t) = -\nabla \cdot j(x) =  -\nabla \cdot [\nabla \mu(x)]
\end{equation}
Plugging in the formula for $\nabla \mu$ gives
\begin{equation}
\partial_t p(x,t)  = \frac{1}{\gamma} \nabla \cdot [\nabla U(x) p(x,t)] + \frac{k_BT}{\gamma}\nabla^2 p(x,t).
\end{equation}
which we can recognize as the forward equation for the It\^o SDE
\begin{equation}
dX_t = -\frac{1}{\gamma}\nabla U(X_t)dt + \sqrt{\frac{k_B T}{\gamma}}dW_t.
\end{equation}
Let us summarize
\begin{itemize}
\item We have found that a gradient flow on probability densities corresponds to stochastic gradient descent of at the single particle level. 
\item Conversely, if we take a gradient ODE $x' = -\gamma^{-1}\nabla U(x)$ and add white noise, the probability density undergoes deterministic gradient decent in the free energy. 
\end{itemize}


\subsection*{Example 8.3: Particle with inertia}
If we don't make the over-damped assumption, then the FP equation becomes 
\begin{equation}
\partial_t p(x,v,t)  = \frac{1}{\gamma} v \cdot \nabla_x p(x,v,t) -  \gamma \nabla \cdot \left(( v +\gamma^{-1} \nabla U(x))p(x,v,t) \right) + \frac{k_BT}{\gamma}\nabla^2 p(x,v,t).
\end{equation}

\subsection*{Example 8.4: Diffusion on a sphere}

Now we consider another application of probability current, which is to construct the equation for diffusion on the sphere in rectangular coordinates. 
We will construct an SDE such that
\begin{equation}
X_0 \in S^{n-1} = \{x \in \mathbb R^n : \|x\| = 1\}
\quad \implies \quad
X_t \in S^{n-1}.
\end{equation}

There are many such SDEs, but we focus on the simplest corresponding to pure diffusion. 
For the solution to stay on the sphere, it must be the case that the diffusion term $\sigma(x)\,dW_t$ always produces a noise vector tangent to the sphere. 
Thus a natural choice is 
\begin{equation}
\sigma(x) = I - xx^\top,
\end{equation}
since this is the orthogonal projection onto the tangent space.

Indeed, if $\|x\|=1$ and $y\in\mathbb R^n$, then
\begin{equation}
x^\top \sigma(x) y 
= x^\top y - (x^\top x)(x^\top y)
= 0,
\end{equation}
so $\sigma(x)y \in T_x S^{n-1}$, where
\begin{equation}
T_x S^{n-1} = \{ v \in \mathbb R^n : v^\top x = 0 \}.
\end{equation}

Since $\sigma(x)$ is symmetric and idempotent on the sphere, the diffusion matrix is
\begin{equation}
A(x) = \sigma(x)\sigma(x)^\top 
= (I - xx^\top)^2 
= I - xx^\top,
\qquad x \in S^{n-1}.
\end{equation}

Now consider the probability current associated with an It\^o diffusion
\begin{equation}
dX_t = b(X_t)\,dt + \sigma(X_t)\,dW_t,
\end{equation}
whose density $p(x,t)$ satisfies
\begin{equation}
\partial_t p(x,t) = -\nabla \cdot j(x,t),
\end{equation}
with current
\begin{equation}
j_i(x,t)
=
b_i(x)\,p(x,t)
-
\frac12 \sum_{j=1}^n
\partial_{x_j}\!\big( A_{ij}(x)\,p(x,t) \big).
\end{equation}

To ensure that probability does not flow off the sphere, we impose the zero normal flux condition
\begin{equation}
j(x,t)\cdot n(x) = 0,
\qquad n(x)=x,
\qquad x\in S^{n-1}.
\end{equation}



We compute the normal component of the probability current. Recall that
\begin{equation}
j_i(x,t)
=
b_i(x)\,p(x,t)
-
\frac12 \sum_{j=1}^n
\frac{\partial}{\partial x_j}
\big( A_{ij}(x)\,p(x,t) \big).
\end{equation}
Taking the dot product with $x$ gives
\begin{equation}
j(x,t)\cdot x
=
\sum_{i=1}^n x_i b_i(x)\,p(x,t)
-
\frac12
\sum_{i=1}^n
x_i
\sum_{j=1}^n
\frac{\partial}{\partial x_j}
\big( A_{ij}(x)\,p(x,t) \big).
\end{equation}


Note that 
\begin{align}
\frac{\partial}{\partial x_j}
\big( (A(x)x)_j\,p \big)
&=
\sum_{i=1}^n
\frac{\partial}{\partial x_j}
\big( A_{ij}(x)\,x_i\,p \big) \\
&=
\sum_{i=1}^n
\left[
\frac{\partial A_{ij}}{\partial x_j} x_i p
+
A_{ij} \frac{\partial x_i}{\partial x_j} p
+
A_{ij} x_i \frac{\partial p}{\partial x_j}
\right].
\end{align}

On the other hand,
\begin{align}
\sum_{i=1}^n x_i
\sum_{j=1}^n
\frac{\partial}{\partial x_j}
\big( A_{ij} p \big)
&=
\sum_{i=1}^n
\sum_{j=1}^n
x_i
\left(
\frac{\partial A_{ij}}{\partial x_j} p
+
A_{ij} \frac{\partial p}{\partial x_j}
\right).
\end{align}

Comparing the two expressions, the only extra term in the full product rule expansion is
\[
\sum_{i=1}^n
\sum_{j=1}^n
A_{ij}(x)\,
\frac{\partial x_i}{\partial x_j}\,p.
\]

Since
\[
\frac{\partial x_i}{\partial x_j}
=
\begin{cases}
1 & i=j, \\
0 & i\ne j,
\end{cases}
\]
this equals
\begin{equation}
\sum_{i=1}^n A_{ii}(x)\,p
=
\operatorname{tr}(A(x))\,p.
\end{equation}

Therefore
\begin{equation}
\sum_{i=1}^n
x_i
\sum_{j=1}^n
\frac{\partial}{\partial x_j}
\big( A_{ij}(x)\,p \big)
=
\sum_{j=1}^n
\frac{\partial}{\partial x_j}
\big( (A(x)x)_j\,p \big)
-
\operatorname{tr}(A(x))\,p.
\end{equation}

Using the identity
\begin{equation}
x_i \partial_{x_j}(A_{ij}p)
=
\partial_{x_j}\!\big( (A(x)x)_j\,p \big)
-
\operatorname{tr}(A(x))\,p,
\end{equation}
and observing that $A(x)x = 0$ on $S^{n-1}$, we obtain
\begin{equation}
x_i \partial_{x_j}(A_{ij}p)
=
- \operatorname{tr}(A(x))\,p.
\end{equation}

Since
\begin{equation}
\operatorname{tr}(A(x))
=
\operatorname{tr}(I - xx^\top)
=
n - \|x\|^2
=
n-1
\quad \text{on } S^{n-1},
\end{equation}
we conclude
\begin{equation}
j(x,t)\cdot x
=
\Big(
b(x)\cdot x
+
\frac{n-1}{2}
\Big)
p(x,t).
\end{equation}

The zero-flux condition must hold for all $p$, hence
\begin{equation}
b(x)\cdot x
=
-\frac{n-1}{2},
\qquad x\in S^{n-1}.
\end{equation}

By rotational symmetry, the simplest choice is a purely radial drift,
\begin{equation}
b(x)
=
-\frac{n-1}{2}\,x.
\end{equation}

Therefore the It\^o SDE for diffusion on the sphere in rectangular coordinates is
\begin{equation}
dX_t
=
-\frac{n-1}{2}\,X_t\,dt
+
\bigl(I - X_t X_t^\top\bigr)\,dW_t,
\qquad
X_0 \in S^{n-1}.
\end{equation}


\section*{Probability current and boundary conditions (Section 8.2)}
Now consider a general Fokker-Planck equation 
\begin{equation}
\partial_t p(x,t) = -\nabla \cdot [b(x)p(x,t)] + \frac{1}{2} \nabla^2: A(x)p(x,t)
\end{equation}
where $A(x) = \sigma(x)\sigma(x)^T$. 
which we can rewrite as
\begin{equation}
\partial_t p(x,t) = -\nabla \cdot j(x,t),\quad j(x,t) = - b(x)p(x,t) + \frac{1}{2} \nabla  \sigma^2(x)p(x,t)
\end{equation}
Thus, $j(x,t) \cdot n$ is a probability in the direction $n$. Just like the probability currents we defined in the context of $Q$-processes, it gives us the rate at which probability flows from position $x$ in the direction $n$. In the previous example we saw that for Brownian dynamics in equilibrium $j(x,t) =0$ because the chemical potential vanishes. In general, this need not be the case and we can have non-equilibrium steady-states with non-vanishing probability currents. This will always be the case when $b(x)$ is not the gradient of a potential. 

\subsection*{Diffusion with reflecting boundaries (see also Section 6.6)}

So far we have considered diffusion processes in unbounded domains, or situations where the SDE itself ensures that the process is constrained to a certain region (for example, geometric Brownian motion is constrained to $[0,\infty)$). Often we would like to explicitly enforce the boundary conditions. This can be achieved by placing a constraint on the flux.

For example, suppose we want to consider the diffusion process
\begin{equation}
dX_t = \sigma\, dW_t.
\end{equation}
on the interval $[0,L]$. This means $X_t$ cannot pass the boundary, but this is not enforced in the SDE itself. We can enforce it in the PDE, however. The forward equation is
\begin{equation}
\partial_t p = \frac{1}{2}\sigma^2\,\partial_x^2 p,
\end{equation}
and the probability flux is
\begin{equation}
j(x) = \frac{\sigma^2}{2}\,\partial_x p.
\end{equation}
Reflecting boundary conditions correspond to zero flux, so we impose
\begin{equation}
j(0)=j(L)=0,
\end{equation}
which is equivalent to the Neumann conditions
\begin{equation}
\partial_x p(0,t)=\partial_x p(L,t)=0.
\end{equation}

Now consider the special case where the initial condition is a point mass at $x_0\in(0,L)$,
\begin{equation}
p(x,0) = p_0(x) = \delta(x - x_0).
\end{equation}
In this case, the cosine–series coefficients are
\begin{equation}
a_0 = \frac{1}{L}, \qquad
a_n = \frac{2}{L}\cos\!\left(\frac{n\pi x_0}{L}\right), \quad n \ge 1.
\end{equation}
The solution (the transition density with reflecting boundaries) is then
\begin{equation}
p(x,t\,|\,x_0)
=
\frac{1}{L}
+
\frac{2}{L}
\sum_{n=1}^\infty
\cos\!\left(\frac{n\pi x_0}{L}\right)
\cos\!\left(\frac{n\pi x}{L}\right)
\exp\!\left[-\frac{\sigma^2}{2}\left(\frac{n\pi}{L}\right)^2 t\right].
\end{equation}

\subsection*{Local time}
It is also possible to enforce the reflecting boundary conditions directly in the SDE. The resulting process is reflected Brownian motion on $[0,L]$, which is written
\begin{equation}
dX_t = \sigma\, dW_t + dL_t^{(0)} - dL_t^{(L)},
\end{equation}
where $L_t^{(0)}$ and $L_t^{(L)}$ are boundary local times that increase only when $X_t=0$ or $X_t=L$, respectively, and push the process back into the interval. This SDE is the stochastic analogue of imposing zero–flux (Neumann) boundary conditions in the PDE.


\section*{Equilibrium steady-states}


Now consider the infinitesimal generator of Brownian dynamics
\begin{equation}\label{eq:brownian_generator}
\mathcal L f = \gamma^{-1} \nabla U \cdot \nabla f + D \nabla^2 f.
\end{equation}
where $D = k_BT/\gamma$. As we have already seen, the process processes an equilibrium steady-state where the current $j(x)=0$ for all $x \in {\mathbb R}^d$. We saw with $Q$-processes that equilibrium steady-states are precisely the ones where the reversed process looks the same as the forward process. Let's first consider the time-reversal of this diffusion process. 


\subsection*{Time-reversal}
Consider a general It\^o process
\begin{equation}
dX_t = b(X_t)dt + \sigma(X_t)dW_t
\end{equation}
which processes a unique time-invariant  steady-state $\pi$ and suppose that $X_0 \sim \pi$. We want to consider the time reversed process
\begin{equation}
\tilde{X}_t = X_{T-t}. 
\end{equation}
Over an increment $dt$,
\begin{equation}
(\tilde{X}_{t+dt} - \tilde{X}_t)|\tilde{X}_t  = (X_{T-t-dt} - X_{T-t})|X_{T-t} 
\end{equation}
Using Bayes' formula 
\begin{equation}
(X_{T-t-dt} - X_{T-t}) \mid X_{T-t}=x
\sim 
\mathcal N\big(-b(x)\,dt + A(x)\,\nabla \log\pi(x)\,dt,\; A(x)\,dt\big).
\end{equation}
where $A(x) = \sigma(x)\sigma(x)^T$. This suggests the time-reversal is 
\begin{equation}
d\tilde{X}_t = -b(x)dt + A(x)\nabla \log\pi(x)dt + \sigma(x)d\tilde{W}_t
\end{equation}
where $\tilde{W}_t$ is another Wiener process. 

Another way to derive this is from the generator view. We have the identity 
\begin{equation}
{\mathbb E}_{\pi}[f(X_0)g(X_t)] = {\mathbb E}_{\pi}[f(\tilde{X}_t)g(\tilde{X}_0)] 
\end{equation}
By the backward equation 
\begin{equation}
\begin{aligned}
\left.\frac{\partial}{\partial t}\,\mathbb{E}_{\pi}\big[f(X_0)g(X_t)\big]\right|_{t=0}
&=
\left.\mathbb{E}_{\pi}\!\left[ \frac{\partial}{\partial t}\,\mathbb{E}\big[f(X_0)g(X_t)\mid X_0\big]\right]\right|_{t=0} \\
&=
\mathbb{E}_{\pi}\big[f(X_0)\,\mathcal{L}g(X_0)\big]
=
\langle f, \mathcal{L}g\rangle_{\pi}.
\end{aligned}
\end{equation}
where $\langle \cdot,\dot \rangle_{\pi}$ is the inner product on the Hilbert space
\begin{equation}
L^2(\pi) = \{f: \int f(x)\pi(x)dx < \infty \}.
\end{equation}
That is, 
\begin{equation}
\langle f,g \rangle_{\pi} = \int f(x)g(x)\pi(x)dx.
\end{equation}
Therefore the generator  $\tilde{ \mathcal{L}}$ of $\tilde{Y}$ must satisfy 
\begin{equation}
\langle f, \mathcal{L}g\rangle_{\pi} = \langle \tilde{ \mathcal{L}}f,g\rangle_{\pi}
\end{equation}
and hence $\tilde{ \mathcal{L}} =\mathcal{L}^{\dagger}$ is the adjoint with respect to $\pi$. 
In particular, 
\begin{align}
\langle f,\mathcal L g\rangle_{\pi}
&= \int_{\mathbb R^d} f(x)\,\big(b(x)\cdot\nabla g(x) + D\,\Delta g(x)\big)\,\pi(x)\,dx \\[0.3em]
&= -\int_{\mathbb R^d} g(x)\,\nabla\!\cdot\!\big(f(x)\,b(x)\,\pi(x)\big)\,dx
   \;+\;D\int_{\mathbb R^d} g(x)\,\Delta\big(f(x)\,\pi(x)\big)\,dx \\[0.3em]
&= \int_{\mathbb R^d} g(x)\Big(
   -b(x)\cdot\nabla f(x)\,\pi(x)
   -f(x)\,\nabla\!\cdot\!\big(b(x)\,\pi(x)\big) \\
&\hspace{5em}
   \qquad\qquad
   +D\big(\Delta f(x)\,\pi(x)
         +2\,\nabla f(x)\cdot\nabla\pi(x)
         +f(x)\,\Delta\pi(x)\big)
   \Big)\,dx \\[0.3em]
&\overset{\nabla\cdot(b\pi)=D\Delta\pi}{=}
   \int_{\mathbb R^d} g(x)\Big(
   -b(x)\cdot\nabla f(x)\,\pi(x)
   +D\,\Delta f(x)\,\pi(x)
   +2D\,\nabla f(x)\cdot\nabla\pi(x)
   \Big)\,dx \\[0.3em]
&= \int_{\mathbb R^d} g(x)\Big(
   \big(-b(x)+2D\,\nabla\log\pi(x)\big)\cdot\nabla f(x)
   +D\,\Delta f(x)
   \Big)\,\pi(x)\,dx \\[0.3em]
&= \langle \mathcal L^{\dagger} f, g\rangle_{\pi},
\end{align}

\subsection*{Symmetrization and Schr\"odinger representation}

In the equilibrium setting, the Brownian dynamics associated with a potential $U:\mathbb R^d\to\mathbb R$ is reversible with respect to the Gibbs measure
\begin{equation}
\pi(x)
=
Z^{-1}\exp\!\Big(-\frac{U(x)}{k_B T}\Big),
\end{equation}
where $Z$ is the normalizing constant. In the overdamped Langevin form
\begin{equation}
dX_t = -D\,\nabla U(X_t)\,dt + \sqrt{2D}\,dW_t,
\end{equation}
with $D = k_B T/\gamma$, the generator acting on observables is
\begin{equation}
\mathcal L f(x)
=
D\,\Delta f(x) - D\,\nabla U(x)\cdot\nabla f(x).
\end{equation}
The invariant density $\pi$ satisfies detailed balance, and one checks that
\begin{equation}
\langle f,\mathcal L g\rangle_{\pi}
=
\langle \mathcal L f,g\rangle_{\pi},
\qquad
\langle f,g\rangle_{\pi}
=
\int_{\mathbb R^d} f(x)\,g(x)\,\pi(x)\,dx,
\end{equation}
so $\mathcal L$ is a symmetric (Hermitian) operator on the Hilbert space $L^2(\pi)$.

It is often convenient to conjugate $\mathcal L$ to a Schr\"odinger operator on $L^2(dx)$ with respect to Lebesgue measure. Define the unitary map $\mathcal U : L^2(\pi)\to L^2(dx)$ by
\begin{equation}
(\mathcal U f)(x)
=
\pi(x)^{1/2} f(x)
=
Z^{-1/2}\exp\!\Big(-\frac{U(x)}{2k_B T}\Big) f(x).
\end{equation}
Unitarity follows from
\begin{equation}
\int_{\mathbb R^d} |(\mathcal U f)(x)|^2\,dx
=
\int_{\mathbb R^d} f(x)^2\,\pi(x)\,dx
=
\|f\|_{L^2(\pi)}^2.
\end{equation}
We then define the Schr\"odinger operator $\mathcal H$ by the similarity transform
\begin{equation}
\mathcal H
=
-\,\mathcal U\,\mathcal L\,\mathcal U^{-1},
\end{equation}
so that $\mathcal H$ is self-adjoint on $L^2(dx)$ whenever $\mathcal L$ is self-adjoint on $L^2(\pi)$.

To compute $\mathcal H$ explicitly, set $\psi = \mathcal U f$, so that
\begin{equation}
f(x) = (\mathcal U^{-1}\psi)(x)
=
\pi(x)^{-1/2}\psi(x)
=
Z^{1/2}\exp\!\Big(\frac{U(x)}{2k_B T}\Big)\psi(x).
\end{equation}
Then
\begin{equation}
\mathcal H\psi(x)
=
-\,\pi(x)^{1/2}\,\mathcal L\big(\pi^{-1/2}\psi\big)(x).
\end{equation}
Write $\beta = (k_B T)^{-1}$ and note that $\pi^{-1/2} = Z^{1/2}e^{\beta U/2}$. A direct computation gives
\begin{align}
\nabla\big(\pi^{-1/2}\psi\big)
&=
e^{\beta U/2}\Big(\nabla\psi + \tfrac{\beta}{2}\psi\,\nabla U\Big), \\[0.3em]
\Delta\big(\pi^{-1/2}\psi\big)
&=
e^{\beta U/2}\Big(
\Delta\psi
+
\beta\,\nabla U\cdot\nabla\psi
+
\tfrac{\beta}{2}(\Delta U)\,\psi
+
\tfrac{\beta^2}{4}|\nabla U|^2 \psi
\Big).
\end{align}
Using $\mathcal L f = D\Delta f - D\,\nabla U\cdot\nabla f$, we obtain
\begin{align}
\mathcal L\big(\pi^{-1/2}\psi\big)
&=
D\,\Delta\big(\pi^{-1/2}\psi\big)
-
D\,\nabla U\cdot\nabla\big(\pi^{-1/2}\psi\big) \\[0.3em]
&=
D\,e^{\beta U/2}\Big(
\Delta\psi
+
\beta\,\nabla U\cdot\nabla\psi
+
\tfrac{\beta}{2}(\Delta U)\,\psi
+
\tfrac{\beta^2}{4}|\nabla U|^2 \psi
\Big) \\[0.3em]
&\quad
- D\,e^{\beta U/2}\Big(
\nabla U\cdot\nabla\psi
+
\tfrac{\beta}{2}|\nabla U|^2\psi
\Big) \\[0.3em]
&=
D\,e^{\beta U/2}\Big(
\Delta\psi
+
\Big(\beta-1\Big)\nabla U\cdot\nabla\psi
+
\tfrac{\beta}{2}(\Delta U)\,\psi
+
\Big(\tfrac{\beta^2}{4}-\tfrac{\beta}{2}\Big)|\nabla U|^2 \psi
\Big).
\end{align}
In the equilibrium Langevin case we have $\beta = 1/(k_B T)$ and $D = k_B T/\gamma$, so that $\beta D = 1/\gamma$ and $\beta^2 D = 1/(\gamma k_B T)$. Substituting $\pi^{1/2} = Z^{-1/2}e^{-\beta U/2}$, the prefactors cancel and we obtain
\begin{align}
\mathcal H\psi(x)
&=
-\,\pi(x)^{1/2}\,\mathcal L\big(\pi^{-1/2}\psi\big)(x) \\[0.3em]
&=
- D\,\Delta\psi(x)
+
\Big[
\frac{1}{4\gamma k_B T}|\nabla U(x)|^2
-
\frac{1}{2\gamma}\,\Delta U(x)
\Big]\psi(x).
\end{align}
Thus $\mathcal H$ has the Schr\"odinger form
\begin{equation}
\mathcal H
=
- D\,\Delta + V(x),
\end{equation}
with effective potential
\begin{equation}
V(x)
=
\frac{1}{4\gamma k_B T}|\nabla U(x)|^2
-
\frac{1}{2\gamma}\,\Delta U(x).
\end{equation}

Finally, if $f_t$ solves the backward Kolmogorov equation
\begin{equation}
\partial_t f_t
=
\mathcal L f_t,
\end{equation}
and we define
\begin{equation}
\psi_t(x)
=
(\mathcal U f_t)(x)
=
\pi(x)^{1/2} f_t(x),
\end{equation}
then
\begin{equation}
\partial_t \psi_t
=
\mathcal U\,\partial_t f_t
=
\mathcal U\,\mathcal L f_t
=
-\,\mathcal H(\mathcal U f_t)
=
-\,\mathcal H\psi_t.
\end{equation}
Therefore $\psi_t$ satisfies the imaginary-time Schr\"odinger equation
\begin{equation}
\partial_t \psi_t(x)
=
-\,\mathcal H\,\psi_t(x),
\end{equation}
and the spectral properties of the symmetric generator $\mathcal L$ on $L^2(\pi)$ are encoded by the spectrum of the Hermitian Schr\"odinger operator $\mathcal H$ on $L^2(dx)$.











\end{document}